\subsection{Sharpness and restricted adversaries}\label{sec:consequences}

\begin{theorem}[The attaining controller]\label{thm:attaining}
The controller $\Astar$ succeeds from every state-$0$ start pair on every legal induced grid maze with at most six cells. Moreover,
\[
 \tau(\Astar)=\tau_{\mathcal T_{\rm grid}}(\Astar)
   =\tau_{\mathcal J}(\Astar)=7.
\]
\end{theorem}
\begin{proof}
Proposition~\ref{c:six-cell} proves success by combining the priority classification of Lemma~\ref{c:geometry}, Proposition~\ref{t:small-trees}, and Lemmas~\ref{c:full-centre}, \ref{c:square}, and~\ref{c:rectangles}. Lemma~\ref{a:trap} supplies a seven-cell path trap, which belongs to each displayed restricted class. There is no smaller trap even in the unrestricted class.
\end{proof}

Proposition~\ref{u:upper} gives the opposite bound for every controller with at most two states, with all its witnesses in $\mathcal J$. The attaining controller therefore proves Theorem~\ref{thm:B}. This deduction uses no auxiliary classification of controllers on paths.

The constant seven has a geometric explanation within the proof. The two vertical state-$1$ edge cycles of the U witness each have only a northern external port. The port comparisons in the tree proof exclude a connection of length at most three between two such separated supports. A connector of length four can have direction word $\N\E\E\Sdir$, using three new cells in addition to the four support cells. At six cells, the mixed profiles that fit either have a state-$0$ gate blocking one basin or produce the exceptional path $P_H$, where basin inversion forces contact before recurrence. Seven is therefore the first available size for this separated-shuttle mechanism.

This explanation does not classify all minimal traps. A spatially cyclic seven-cell maze can fail by a different mechanism. For example, on
\[
 \{(0,1),(1,0),(1,1),(1,2),(2,0),(2,1),(3,1)\}
\]
the four tags
\[
 (1,0)_0\longrightarrow(1,1)_0\longrightarrow(2,1)_0
 \longrightarrow(2,0)_1\longrightarrow(1,0)_0
\]
form a recurrent tagged cycle. Starting at its two state-$0$ opposite corners $(1,0)$ and $(2,1)$ maintains Manhattan distance two. This also shows why the same-cycle contact lemma cannot be used without its path-support hypothesis.
